\section{Limitations}
\label{sec:limitations}

\paragraph{Synthetic data.}
All 200 patients are generated by Synthea. Synthea produces medically realistic
FHIR records, with medication names, dosages, RxNorm codes, and patient trajectories
following real clinical distributions, but it is not real patient data. Real EHR records
contain transcription errors, inconsistent formatting, missing fields, local coding
conventions, and data quality issues not present in Synthea output. We expect our F1
values to represent upper-bound estimates of performance on real hospital data for the
same task. External validation on real EHR data is the most important next step for
this work.

\paragraph{Exact string matching.}
We evaluate using exact string match against medication names as they appear in the
FHIR record. A prediction of ``metformin 500mg'' against a ground truth of ``Metformin
500 MG Oral Tablet'' counts as a miss. In real-world deployment, medication name
normalisation or fuzzy matching (e.g., RxNorm concept-based matching) would recover
many of these near-misses. Our F1 scores are therefore conservative lower bounds.
The rank ordering of models and strategies is unlikely to change under fuzzy evaluation,
but absolute values would improve.

\paragraph{English only.}
Synthea generates English-language patient records. Clinical NLP performance for
FHIR medication reconciliation in other languages, or for multilingual records, is
not addressed by this study.

\paragraph{Instruction-tuned models only (except BioMistral).}
All models tested other than BioMistral are instruction-tuned. No instruction-tuned
biomedical models (MedAlpaca, Meditron-70B, BioLlama) were included in this evaluation.
BioMistral's failure is a completion-model failure, not a biomedical-LLM failure in
general. Future work should test properly instruction-tuned medical models to assess
whether domain pretraining provides a benefit when paired with instruction fine-tuning.

\paragraph{Strategy A input cap.}
Strategy A applies an active-first, 100-entry cap on \texttt{MedicationRequest}
resources (see Section~\ref{sec:strategies}); Strategies B, C, and D do not.
This asymmetry is necessary to fit Synthea's refill-level bundles into the 32K
context windows of the 7B models, but it means Strategy A receives a mildly
preprocessed JSON rather than the raw bundle. The cap is conservative for our
comparison: it guarantees every active medication enters the model input, so any
remaining recall gap against the other strategies reflects format effects, not
missing evidence. A deployment that wants true raw JSON on a 32K-context model
would either truncate at the token level or reject long-history patients, both of
which would degrade Strategy A further.

\paragraph{Single prompt template.}
We use one prompt design across all models and strategies. We do not ablate prompt
wording, system message length, instruction phrasing, or output format constraints.
Strategy effects could interact with prompt effects in ways this study cannot disentangle.
A prompt-ablation study on a single model would help separate these contributions.
